Put $E=(0,0),B=(0,h),D=(d,0),A=(a,h),C=(c,0)$. Since $F$ is the centroid of $ADE$ and lies on $BD$, comparing its coordinates gives $a=d$ and $F=(\dfrac{2d}{3},\dfrac{h}{3})$. The perpendicularity $EF\perp BD$ yields $h^2=2d^2$. Perpendicular diagonals give $d(c-d)+h^2=0$, so $c=-d$. Now $AD=h=10$ and $|d|=5\sqrt2$. The bases have lengths $5\sqrt2,10\sqrt2$, giving area $\boxed{75\sqrt2}$. 
\begin{asy}
texpreamble("\usepackage{amsmath,amssymb}");
size(220);
defaultpen(fontsize(9pt));
real d=5*sqrt(2);pair A=(d,10),B=(0,10),C=(-d,0),D=(d,0),E=(0,0),F=(2*d/3,10/3);draw(A--B--C--D--cycle);draw(A--C);draw(B--D);draw(B--E--F);draw(A--E,dashed);pair[] P={A,B,C,D,E,F};string[] L={"A","B","C","D","E","F"};for(int i=0;i<6;++i)label("$"+L[i]+"$",P[i],i==5?(1,0):(i<2?(0,1):(0,-1)));
\end{asy}
